\section{Resultados e Discussão}
\label{sec:discussao}

Esta seção apresenta as estimativas econométricas que confrontam as hipóteses formuladas sobre os determinantes da conversão da vegetação nativa em Mato Grosso entre 2001 e 2024. A discussão organiza-se em seis partes integradas:
\begin{enumerate}[label=(\arabic*),leftmargin=1.8cm,itemsep=2pt,topsep=2pt]
    \item Avaliação da hipótese de mercado (H1a), examinando como os choques de preços internacionais interagem com a aptidão biofísica do solo (Tabela~\ref{tab:h1-principais});
    \item Condicionamento tecnológico, ocupação agrícola pré-determinada e evidências de intensificação produtiva em terras já abertas (Tabela~\ref{tab:h1-margens-tecnologia});
    \item Custos logísticos de transporte e o papel moderador do freio regulatório exercido por embargos ambientais (H1b) (Tabela~\ref{tab:h1-logistica-embargos});
    \item Quebras estruturais por regimes institucionais e assimetrias regulatórias entre a Amazônia e o Cerrado (Tabela~\ref{tab:h1-regimes-biomas});
    \item Assimetria setorial entre agricultura e pecuária e a bateria de testes de robustez e sensibilidade inferencial (Tabela~\ref{tab:h1-destinos-robustez});
    \item Diálogo com a literatura empírica e articulação com a arquitetura geral da tese de doutorado.
\end{enumerate}

\subsection{Preços reais, aptidão biofísica e conversão para agricultura (H1a)}
\label{discussao-principais}

A amostra de referência para a estimação do modelo base da Equação~\eqref{eq:modelo-h1a} compreende $2.767$ observações município--ano, cobrindo o período de 2001 a 2024. Essa especificação exclui os municípios sem cobertura florestal inicial mapeada pelo MapBiomas ou com séries incompletas, além de descartar os grupos unitários decorrentes da inclusão dos efeitos fixos bidirecionais de município ($\alpha_m$) e de bioma por ano ($\lambda_{b(m),t}$). A variável dependente $Y_{mt}^{\text{agr}}$ é expressa em pontos percentuais da vegetação nativa inicial do município convertida anualmente para lavouras agrícolas. Os erros-padrão são agrupados em nível de município ($G = 137$), assegurando robustez a heterocedasticidade e dependência serial temporal arbitrária \cite{CameronMiller2015}.

A Tabela~\ref{tab:h1-principais} reporta as estimativas da especificação principal. As colunas (1) e (2) avaliam individualmente os efeitos das exposições de soja (M1) e milho (M2); a coluna (3) apresenta a especificação conjunta com ambas as culturas (M3); e as colunas (4) a (7) desagregam a amostra de acordo com a predominância territorial dos biomas Amazônia e Cerrado.

\begin{table}[htbp]
\centering
\scriptsize
\setlength{\tabcolsep}{2.8pt}
\captionstyle{\centering}
\normalcaption
\caption{Associação entre preços reais, capacidade potencial e conversão para agricultura, 2001--2024}
\label{tab:h1-principais}
\begin{tabular}{lccccccc}
\hline
 & (1) & (2) & (3) & (4) & (5) & (6) & (7) \\
Variável & M1 Soja & M2 Milho & M3 Conj. & Amaz. (S) & Cerr. (S) & Amaz. (M) & Cerr. (M) \\
\hline
$\ln P_{s,t-1} \times Z_{i,s}$ & 0,0848** & -- & 0,0815* & 0,0866 & 0,0752* & -- & -- \\
 & (0,0358) &  & (0,0433) & (0,0610) & (0,0439) &  &  \\
$\ln P_{m,t-1} \times Z_{i,m}$ & -- & 0,0519*** & 0,0060 & -- & -- & 0,0185 & 0,0572*** \\
 &  & (0,0179) & (0,0239) &  &  & (0,0514) & (0,0192) \\
$Z_{i,s} \times \tau_t$ & 0,0405*** & -- & 0,0247 & 0,0396* & 0,0473** & -- & -- \\
 & (0,0144) &  & (0,0198) & (0,0232) & (0,0183) &  &  \\
$Z_{i,m} \times \tau_t$ & -- & 0,0538*** & 0,0346 & -- & -- & 0,0426 & 0,0623*** \\
 &  & (0,0152) & (0,0220) &  &  & (0,0361) & (0,0175) \\
Precipitação ($R_{it}$) & -0,0135** & -0,0152** & -0,0137** & -0,0230** & -0,0023 & -0,0241** & -0,0048 \\
 & (0,0063) & (0,0067) & (0,0062) & (0,0104) & (0,0058) & (0,0108) & (0,0073) \\
\hline
Observações & 2.767 & 2.767 & 2.767 & 1.303 & 1.412 & 1.303 & 1.412 \\
$R^2$ total & 0,3828 & 0,3807 & 0,3841 & 0,3735 & 0,3924 & 0,3700 & 0,3968 \\
$R^2$ \textit{within} & 0,0198 & 0,0165 & 0,0218 & 0,0149 & 0,0492 & 0,0094 & 0,0560 \\
\hline
\end{tabular}
\par\smallskip\begin{minipage}{\textwidth}\footnotesize
\centering
Fonte: elaboração própria com MapBiomas 10.1, Banco Mundial, BLS/FRED, FAO/IIASA (GAEZ v4) e IBGE.
\end{minipage}
\end{table}

\FloatBarrier

Nos modelos de referência, a interação entre o preço real internacional defasado ($\ln P_{jt-1}$) e o rendimento potencial sob manejo de alto insumo ($Z_{jm}$) é positiva e estatisticamente significativa tanto para a soja ($\hat{\beta}_s = 0{,}0848$, $p = 0{,}019$, coluna 1) quanto para o milho ($\hat{\beta}_m = 0{,}0519$, $p = 0{,}004$, coluna 2). Esse resultado oferece suporte empírico à hipótese H1a: em resposta a elevações nas cotações globais das commodities, a conversão líquida da vegetação nativa para lavouras ocorre de forma sistematicamente mais intensa nos municípios dotados de solos de maior produtividade potencial.

Para dimensionar a magnitude econômica desse efeito, comparam-se dois municípios hipotéticos que se diferenciem em um desvio-padrão de rendimento biofísico potencial, tomando como parâmetro a média estadual e uma localidade situada um desvio-padrão acima dela ($\approx 504\text{ kg/ha}$ de matéria seca na soja e $\approx 852\text{ kg/ha}$ no milho). Uma elevação de 50\% no preço real da commodity ($\Delta \ln P \approx 0{,}405$) gera uma diferença de $0{,}034$ ponto percentual na taxa anual de conversão no caso da soja e de $0{,}021$ ponto percentual no milho. Diante de uma taxa média amostral de conversão de $0{,}17$ ponto percentual ao ano em Mato Grosso, esses diferenciais de inclinação equivalem a uma variação entre 12\% e 20\% da média observada. Trata-se de uma resposta economicamente substantiva, que converge conceitualmente com a evidência de resposta heterogênea ao longo do espaço biofísico documentada por \citeonline{BermanEtAl2023}, ressalvadas as distinções metodológicas relativas à desagregação por usos de destino.

A inclusão da tendência linear diferenciada por aptidão ($Z_{jm} \times \tau_t$) desempenha papel indispensável na identificação econométrica. Os coeficientes estimados para esse termo são positivos e altamente significativos ($0{,}0405$, $p = 0{,}006$ para soja; $0{,}0538$, $p = 0{,}001$ para milho), capturando uma trajetória secular de maior expansão agrícola nos solos mecanizáveis ao longo das duas décadas. Caso essas tendências não fossem explicitamente controladas, as estimativas de $\hat{\beta}_s$ e $\hat{\beta}_m$ subiriam artificialmente para $0{,}1421$ e $0{,}1112$, absorvendo espuriamente a tendência de modernização tecnológica do Centro-Oeste. Paralelamente, o controle de anomalia de precipitação ($R_{mt}$) apresenta coeficientes negativos ($-0{,}0135$ e $-0{,}0152$, $p < 0{,}05$), evidenciando que anomalias positivas de chuva reduzem a taxa de conversão observada, fato que reflete tanto restrições operacionais ao desmatamento no campo quanto maior cobertura de nuvens que prejudica a detecção orbital via satélite \cite{ChomitzThomas2003,AssuncaoEtAl2023}.

Na especificação conjunta da coluna (3), a significância do complexo de grãos se mantém, mas a alta correlação contemporânea entre as séries internacionais de preços da soja e do milho ($r = 0{,}948$) e entre seus potenciais produtivos ($r = 0{,}591$) restringe a precisão individual dos parâmetros: o coeficiente da soja mantém-se estável ($\hat{\beta}_s = 0{,}0815$, $p = 0{,}062$), enquanto o do milho torna-se estatisticamente indistinguível de zero ($0{,}0060$, $p = 0{,}802$). O teste de Wald de significância conjunta, entretanto, rejeita enfaticamente a hipótese nula de que ambas as interações sejam conjuntamente nulas ($F = 4{,}12$, $p = 0{,}0183$), corroborando a presença de um vetor econômico comum impulsionador do uso da terra.

A desagregação por bioma nas colunas (4) a (7) aponta heterogeneidade no ajuste e na precisão estatística. No Cerrado, onde o relevo tabular e a fertilidade corrigida permitiram a rápida consolidação do agronegócio mecanizado, os coeficientes de soja ($0{,}0752$, $p = 0{,}092$) e milho ($0{,}0572$, $p = 0{,}004$) apresentam maior significância estatística, com $R^2$ \textit{within} alcançando $0{,}0492$ e $0{,}0560$ (mais do que o triplo do ajuste observado na Amazônia). Nos municípios de predominância amazônica, o coeficiente da soja é de $0{,}0866$, porém estimado com menor precisão ($p = 0{,}160$). Testes formais de igualdade de inclinação baseados em interações contínuas pela fração municipal do bioma não rejeitam a hipótese nula de coeficientes idênticos entre Amazônia e Cerrado ($p = 0{,}703$ para soja e $p = 0{,}911$ para milho), indicando que as disparidades observadas refletem variabilidade amostral e dispersão dos resíduos, e não uma diferença estatística estrutural na elasticidade de longo prazo.

\subsection{Condicionamento tecnológico, histórico de ocupação e intensificação produtiva}
\label{discussao-margens}

O modelo de referência ancora a aptidão biofísica no cenário de alto insumo tecnológico do GAEZ v4, representativo da agricultura moderna de precisão. Para examinar se a sensibilidade aos preços depende desse pacote tecnológico e distinguir a abertura de novas fronteiras do aproveitamento de terras já desmatadas, a Tabela~\ref{tab:h1-margens-tecnologia} investiga medidas alternativas de capacidade agronômica e o histórico de ocupação agrícola pré-determinada.

\begin{table}[htbp]
\centering
\scriptsize
\setlength{\tabcolsep}{2.8pt}
\captionstyle{\centering}
\normalcaption
\caption{Cenários de capacidade tecnológica, ocupação agrícola pré-determinada (1985--2000) e área de milho}
\label{tab:h1-margens-tecnologia}
\begin{tabular}{lccccccc}
\hline
 & (1) & (2) & (3) & (4) & (5) & (6) & (7) \\
Variável & Baixo Ins. & Adequação & PAM 91--00 & Agric. 85--00 & Soja 2000 & GAEZ+Agric. & Milho (\%) \\
\hline
$\ln P_{c,t-1} \times Z_{i,c}$ & -0,0191 & -0,0874** & -0,0644 & -- & -- & 0,0496* & -0,1859 \\
 & (0,0286) & (0,0354) & (0,0488) &  &  & (0,0269) & (0,1772) \\
$\ln P_{c,t-1} \times H_i$ & -- & -- & -- & -0,1160*** & -0,1465*** & -0,1038** & 0,5731* \\
 &  &  &  & (0,0429) & (0,0431) & (0,0404) & (0,3156) \\
$Z_{i,c} \times \tau_t$ & 0,0049 & -0,0337** & -0,0427* & -- & -- & 0,0188 & -0,6453* \\
 & (0,0125) & (0,0146) & (0,0218) &  &  & (0,0137) & (0,3276) \\
$H_i \times \tau_t$ & -- & -- & -- & -0,0695*** & -0,0816*** & -0,0654*** & 2,5136*** \\
 &  &  &  & (0,0224) & (0,0199) & (0,0234) & (0,5835) \\
Precipitação ($R_{it}$) & -0,0167** & -0,0157** & -0,0305*** & -0,0150** & -0,0135** & -0,0136** & 0,4372*** \\
 & (0,0068) & (0,0064) & (0,0106) & (0,0061) & (0,0056) & (0,0059) & (0,1385) \\
\hline
Observações & 2.767 & 2.767 & 1.016 & 2.767 & 2.767 & 2.767 & 2.767 \\
$R^2$ total & 0,3742 & 0,3836 & 0,5012 & 0,4060 & 0,4179 & 0,4081 & 0,8618 \\
$R^2$ \textit{within} & 0,0046 & 0,0235 & 0,0398 & 0,0566 & 0,0754 & 0,0612 & 0,3241 \\
\hline
\end{tabular}
\par\smallskip\begin{minipage}{\textwidth}\footnotesize
\centering
Fonte: elaboração própria com MapBiomas 10.1, FAO/IIASA (GAEZ v4), IBGE/PAM e Banco Mundial.
\end{minipage}
\end{table}

\FloatBarrier

As colunas (1) a (3) da Tabela~\ref{tab:h1-margens-tecnologia} mostram que a relação positiva identificada na referência desaparece quando se utilizam métricas que desconsideram o aporte intensivo de insumos. Sob o cenário de baixo insumo, o coeficiente da interação de preço torna-se ligeiramente negativo e estatisticamente não significativo ($-0{,}0191$, $p = 0{,}507$, coluna 1). Quando se utiliza o índice genérico de adequação agroclimática, o coeficiente torna-se negativo ($-0{,}0874$, $p = 0{,}015$, coluna 2), e o rendimento comercial histórico da Pesquisa Agrícola Municipal (PAM) entre 1991 e 2000 também gera estimativa negativa e imprecisa ($-0{,}0644$, $p = 0{,}194$, coluna 3). Esses contrastes confirmam que o mecanismo econômico de H1a está estritamente condicionado ao modelo tecnológico do Centro-Oeste: a resposta da conversão aos preços internacionais não decorre da fertilidade natural pura da terra, mas da sua capacidade de resposta à mecanização e correção química sob gestão empresarial \cite{Spera2017,LubowskiEtAl2008}.

As colunas (4) a (6) introduzem a fração de ocupação agrícola pré-determinada no período 1985--2000 ($H_m$). Quando o preço real da soja é interagido com o histórico agrícola prévio, o coeficiente é negativo e fortemente significativo ($\hat{\theta}_s = -0{,}1160$, $p = 0{,}008$, coluna 4), acompanhado por uma tendência diferencial também negativa ($-0{,}0695$, $p = 0{,}002$). Esse padrão se repete com a participação da soja em 2000 ($-0{,}1465$, $p < 0{,}001$, coluna 5). Na coluna (6), ao incluir simultaneamente a interação com a aptidão biofísica GAEZ e com a ocupação prévia, ambos os parâmetros preservam sua significância com sinais contrapostos: a aptidão potencial estimula a conversão ($\hat{\beta}_s = 0{,}0496$, $p = 0{,}067$), enquanto a ocupação agrícola anterior arrefece a abertura de novas áreas ($\hat{\theta}_s = -0{,}1038$, $p = 0{,}011$).

Esse descolamento entre aptidão agronômica e presença de lavouras consolidadas é aprofundado na coluna (7), cuja variável dependente é a participação da área plantada de milho no território municipal ($A_{mt}^{\text{milho}}$). Nos municípios com maior histórico agrícola pré-2001, a taxa de expansão da cultura do milho foi superior em $2{,}51$ pontos percentuais por década ($p < 0{,}0001$), e a interação entre preço e área consolidada foi positiva ($\hat{\phi} = 0{,}5731$, $p = 0{,}072$), com $R^2$ \textit{within} de $0{,}3241$.

Em termos econômicos, essa evidência empírica indica que os choques positivos de preços não induziram unicamente a abertura horizontal da fronteira florestal. Em municípios já integrados e com áreas agrícolas desmatadas antes de 2001, os produtores responderam aos estímulos de mercado precipuamente pela via da intensificação produtiva — destacando-se a rápida difusão do sistema de cultivo consorciado com milho de segunda safra (safrinha) após a colheita precoce da soja —, o que dispensou o desmatamento adicional de vegetação nativa \cite{MacedoEtAl2012,VanWeyEtAl2013,Braganca2018}.

\subsection{Logística e modulação institucional por embargos ambientais (H1b)}
\label{discussao-logistica-instituicoes}

A viabilidade econômica da produção agropecuária é condicionada pelos custos de frete até os terminais exportadores e pelo risco de sanções legais associadas à conversão ilegal da terra. A Tabela~\ref{tab:h1-logistica-embargos} avalia essas duas dimensões, introduzindo medidas logísticas nas colunas (1) a (3) e indicadores de dissuasão ambiental nas colunas (4) a (7).

\begin{table}[htbp]
\centering
\scriptsize
\setlength{\tabcolsep}{1.8pt}
\captionstyle{\centering}
\normalcaption
\caption{Custos logísticos de acesso a portos e ferrovias e condicionamento por embargos do IBAMA, 2001--2024}
\label{tab:h1-logistica-embargos}
\begin{tabular}{lccccccc}
\hline
 & (1) & (2) & (3) & (4) & (5) & (6) & (7) \\
Variável & Porto Sul & Log. (S) & Log. (M) & Emb. ha (S) & Emb. ha (M) & Mod. Emb. & Autos (S) \\
\hline
$\ln P_{c,t-1} \times Z_{i,c}$ & 0,1009** & 0,0940** & 0,0656*** & 0,0835** & 0,0483*** & 0,0933** & 0,0494 \\
 & (0,0411) & (0,0394) & (0,0214) & (0,0361) & (0,0172) & (0,0375) & (0,0448) \\
$\ln P_{s,t-1} \times Z_{i,s} \times P_i^{\text{sul}}$ & 0,0486 & -- & -- & -- & -- & -- & -- \\
 & (0,0323) &  &  &  &  &  &  \\
$D_{it}^{\text{min}}$ (porto) & -- & 0,0252*** & 0,0261*** & -- & -- & -- & -- \\
 &  & (0,0079) & (0,0084) &  &  &  &  \\
$F_{it}^{2013}$ (ferrovia) & -- & -0,1689*** & -0,1280*** & -- & -- & -- & -- \\
 &  & (0,0453) & (0,0395) &  &  &  &  \\
$\ln(1+E_{i,t-1}^{\text{ha}})$ & -- & -- & -- & -0,0076*** & -0,0082*** & -- & -- \\
 &  &  &  & (0,0029) & (0,0030) &  &  \\
Embargos/autos ($t-1$) & -- & -- & -- & -- & -- & 0,0082 & 0,0002 \\
 &  &  &  &  &  & (0,0124) & (0,0060) \\
$\ln P_{s,t-1} \times \text{Emb.}$ & -- & -- & -- & -- & -- & -0,0457 & -- \\
 &  &  &  &  &  & (0,0323) &  \\
$Z_{i,s} \times \text{Emb.}$ & -- & -- & -- & -- & -- & 0,0077 & -- \\
 &  &  &  &  &  & (0,0171) &  \\
$\ln P_{s,t-1} \times Z_{i,s} \times \text{Emb.}$ & -- & -- & -- & -- & -- & -0,0314* & -- \\
 &  &  &  &  &  & (0,0179) &  \\
$Z_{i,c} \times \tau_t$ & 0,0405*** & 0,0599*** & 0,0470*** & 0,0396*** & 0,0537*** & 0,0389*** & -- \\
 & (0,0144) & (0,0157) & (0,0139) & (0,0149) & (0,0154) & (0,0139) &  \\
Precipitação ($R_{it}$) & -0,0136** & -0,0149** & -0,0173** & -0,0138** & -0,0154** & -0,0137** & -0,0061 \\
 & (0,0062) & (0,0062) & (0,0069) & (0,0063) & (0,0068) & (0,0062) & (0,0039) \\
\hline
Observações & 2.767 & 2.767 & 2.767 & 2.767 & 2.767 & 2.767 & 1.135 \\
$R^2$ total & 0,3840 & 0,3939 & 0,3873 & 0,3842 & 0,3824 & 0,3834 & 0,4485 \\
$R^2$ \textit{within} & 0,0217 & 0,0374 & 0,0269 & 0,0219 & 0,0191 & 0,0207 & 0,0049 \\
\hline
\end{tabular}
\par\smallskip\begin{minipage}{\textwidth}\footnotesize
\centering
Fonte: elaboração própria com MapBiomas 10.1, IBGE (malha 2023), IBAMA e FAO/IIASA (GAEZ v4).
\end{minipage}
\end{table}

\FloatBarrier

Nas colunas (1) a (3), a incorporação de custos de transporte preserva a estabilidade do parâmetro de exposição de mercado: o coeficiente da soja situa-se entre $0{,}0940$ e $0{,}1009$ ($p < 0{,}05$), e o do milho alcança $0{,}0656$ ($p = 0{,}003$). A expansão da malha logística a partir de 2013, representada pela inauguração do terminal ferroviário da Rumo em Rondonópolis ($F_{mt}^{2013}$), exibe associação negativa expressiva com a conversão de vegetação nativa ($-0{,}1689$ para soja e $-0{,}1280$ para milho, $p < 0{,}002$). Esse sinal negativo é consistente com o fato de que a ferrovia atendeu predominantemente as regiões agrícolas já consolidadas do sul e sudeste de Mato Grosso, reduzindo o custo de oportunidade de intensificar as áreas existentes em detrimento da expansão para novas fronteiras pioneiras \cite{Pfaff1999}.

As colunas (4) a (7) examinam a hipótese H1b, avaliando o papel dissuasório e moderador do aparato estatal de fiscalização. Nas colunas (4) e (5), a inclusão da área acumulada sob embargo ambiental do IBAMA defasada em $t-1$ ($\ln[1 + E_{m,t-1}^{\text{ha}}]$) apresenta um efeito direto negativo e altamente significativo sobre a taxa de conversão para agricultura: $\hat{\rho} = -0{,}0076$ ($p = 0{,}009$) no modelo de soja e $\hat{\rho} = -0{,}0082$ ($p = 0{,}007$) no de milho. Os coeficientes econômicos principais permanecem positivos e estáveis ($\hat{\beta}_s = 0{,}0835$, $p = 0{,}022$; $\hat{\beta}_m = 0{,}0483$, $p = 0{,}006$).

A coluna (6) reporta a especificação crucial da tripla interação formalizada na Equação~\eqref{eq:modelo-h1b}, na qual se interage o estímulo de mercado com o estoque municipal de embargos ambientais. Na ausência de embargos registrados ($\text{Emb} = 0$), o gradiente da resposta aos preços nos municípios de alta aptidão atinge $0{,}0933$ ($p = 0{,}014$). Crucialmente, o parâmetro da tripla interação $\ln P_{st-1} \times Z_{ms} \times \text{Emb}_{mt-1}$ é negativo e estatisticamente significativo a 10\% ($\hat{\delta} = -0{,}0314$, $p = 0{,}081$).

Esse coeficiente de $-0{,}0314$ estabelece que a atuação do IBAMA funciona como um efetivo freio regulatório: a presença de sanções de embargo ambiental amortece a elasticidade da conversão aos preços internacionais em cerca de um terço ($-0{,}0314 / 0{,}0933 \approx -33{,}6\%$). Diferentemente dos autos de infração que representam registros pecuniários de baixa liquidez e longa tramitação jurídica, o embargo administrativo acarreta a exclusão imediata do imóvel das cadeias formais de suprimento e a suspensão do acesso ao crédito rural bancário \cite{AssuncaoEtAl2015,HardingEtAl2021}. Esse resultado confere suporte empírico direto à hipótese H1b, demonstrando que a governança ambiental é capaz de mitigar a transmissão dos ciclos de alta de preços das commodities sobre a destruição de ecossistemas nativos.

Por sua vez, a estimação restrita aos autos de infração de Flora na coluna (7), que reduz a amostra para $1.135$ observações devido à descontinuidade dos relatórios públicos de fiscalização, apresenta coeficiente estatisticamente nulo ($0{,}0002$, $p = 0{,}974$). Essa divergência corrobora a distinção teórica entre contagem de autos (uma medida de esforço operacional estatal contaminada pela própria ocorrência simultânea de infrações no campo) e embargos territoriais efetivos, que incidem com rigor punitivo tangível sobre a exploração econômica da terra \cite{Shimshack2014,Sousa2016}.

\subsection{Quebras estruturais, regimes regulatórios e assimetria de biomas}
\label{discussao-regimes}

Para avaliar as hipóteses de inflexão histórica e examinar se a resposta aos incentivos econômicos variou conforme as transformações do marco regulatório ambiental brasileiro, a Tabela~\ref{tab:h1-regimes-biomas} reporta a estimação da Equação~\eqref{eq:modelo-regimes}, discriminando as elasticidades preço--aptidão entre os três períodos institucionais:
\begin{enumerate}[label=(\alph*),leftmargin=1.5cm,itemsep=2pt]
    \item 2001--2005: Regime de abertura do superciclo, baixa eficácia punitiva e desmatamento recorde;
    \item 2006--2011: Regime de endurecimento regulatório (PPCDAm, Moratória da Soja e restrição ao crédito);
    \item 2012--2024: Regime contemporâneo do Novo Código Florestal e governança cadastral do CAR.
\end{enumerate}

\begin{table}[htbp]
\centering
\scriptsize
\setlength{\tabcolsep}{3.2pt}
\captionstyle{\centering}
\normalcaption
\caption{Associação entre preços reais, capacidade potencial e conversão para agricultura por regime temporal e bioma predominante, 2001--2024}
\label{tab:h1-regimes-biomas}
\begin{tabular}{lcccccc}
\hline
 & (1) & (2) & (3) & (4) & (5) & (6) \\
Variável & Soja MT & Soja Amaz. & Soja Cerr. & Milho MT & Milho Amaz. & Milho Cerr. \\
\hline
$\ln P_{c,t-1} \times Z_{i,c} \times \mathbb{I}[2001\text{--}2005]$ & 0,3541*** & 0,3999 & 0,3392** & 0,4234*** & 0,3667 & 0,4499*** \\
 & (0,1353) & (0,2463) & (0,1634) & (0,1206) & (0,3285) & (0,1337) \\
$\ln P_{c,t-1} \times Z_{i,c} \times \mathbb{I}[2006\text{--}2011]$ & 0,0108 & -0,0643** & 0,0523** & 0,0198 & -0,1143 & 0,0553* \\
 & (0,0211) & (0,0318) & (0,0250) & (0,0266) & (0,0729) & (0,0298) \\
$\ln P_{c,t-1} \times Z_{i,c} \times \mathbb{I}[2012\text{--}2024]$ & 0,0120 & 0,0097 & 0,0057 & -0,0041 & -0,0183 & -0,0022 \\
 & (0,0111) & (0,0153) & (0,0149) & (0,0077) & (0,0302) & (0,0074) \\
Precipitação ($R_{it}$) & -0,0138** & -0,0231** & -0,0039 & -0,0161** & -0,0239** & -0,0076 \\
 & (0,0062) & (0,0104) & (0,0057) & (0,0067) & (0,0107) & (0,0073) \\
\hline
Observações & 2.767 & 1.303 & 1.412 & 2.767 & 1.303 & 1.412 \\
$R^2$ total & 0,3856 & 0,3763 & 0,3956 & 0,3835 & 0,3709 & 0,4056 \\
$R^2$ \textit{within} & 0,0242 & 0,0194 & 0,0541 & 0,0209 & 0,0109 & 0,0697 \\
\hline
\end{tabular}
\par\smallskip\begin{minipage}{\textwidth}\footnotesize
\centering
Fonte: elaboração própria com MapBiomas 10.1, Banco Mundial, BLS/FRED, FAO/IIASA (GAEZ v4) e IBGE.
\end{minipage}
\end{table}

\FloatBarrier

Os resultados do Painel A da Tabela~\ref{tab:h1-regimes-biomas} revelam uma marcante quebra estrutural no comportamento do agronegócio mato-grossense:
\begin{itemize}[leftmargin=1.5cm,itemsep=3pt]
    \item \textbf{Período 2001--2005:} A sensibilidade da conversão aos preços internacionais foi expressiva e altamente significativa em todo o estado. Para a soja, o coeficiente em Mato Grosso foi de $\hat{\beta}_s = 0{,}3541$ ($p < 0{,}01$, coluna 1), atingindo $0{,}3392$ no Cerrado ($p = 0{,}038$, coluna 3) e $0{,}3999$ na Amazônia (coluna 2). No caso do milho, o parâmetro estadual atingiu $0{,}4234$ ($p < 0{,}01$, coluna 4), elevando-se a $0{,}4499$ no Cerrado ($p < 0{,}01$, coluna 6). Nesse primeiro regime, a ausência de amarras de governança permitiu que o mercado operasse com potência máxima sobre a abertura de novas terras;
    \item \textbf{Período 2006--2011:} Observa-se uma nítida divergência espacial entre os biomas. Na Amazônia, o coeficiente de resposta da soja inverte de sinal e torna-se negativo e estatisticamente significativo ($\hat{\beta}_s = -0{,}0643$, $p = 0{,}043$, coluna 2), refletindo o estrangulamento da conversão direta induzido pela conjugação do PPCDAm com o pacto privado da Moratória da Soja (que baniu o comércio de grãos oriundos de desmatamentos amazônicos pós-2006). Em contrapartida, no Cerrado — bioma não contemplado pela Moratória —, a resposta aos preços manteve-se positiva e estatisticamente significativa ($\hat{\beta}_s = 0{,}0523$, $p = 0{,}036$, coluna 3; $\hat{\beta}_m = 0{,}0553$, $p = 0{,}064$, coluna 6);
    \item \textbf{Período 2012--2024:} A elasticidade de conversão da vegetação nativa aos preços perde significância estatística em ambos os biomas e em ambas as culturas ($\hat{\beta} \approx 0{,}012$ na soja e $-0{,}004$ no milho, com $p > 0{,}25$), indicando o arrefecimento generalizado da conversão direta de florestas como vetor primário de expansão de lavouras.
\end{itemize}

O Painel B da Tabela~\ref{tab:h1-regimes-biomas} atesta a robustez dessas conclusões sob uma formulação irrestrita que inclui explicitamente as interações de nível $Z_{jm} \times \text{Regime}$ e centraliza os preços na média de cada período. No Cerrado, a resposta ao preço durante 2006--2011 permaneceu positiva e estatisticamente sólida ($0{,}0569$, $p = 0{,}028$), contrastando com a queda significativa registrada na Amazônia ($-0{,}0704$, $p = 0{,}035$). Essa evidência confirma a eficácia focalizada dos acordos setoriais privados na blindagem do bioma Amazônia, ao mesmo tempo em que alerta para o efeito de vazamento espacial e deslocamento de pressões econômicas em direção ao bioma Cerrado \cite{HeilmayrEtAl2020}.

\subsection{Assimetria setorial entre destinos e consistência inferencial}
\label{discussao-destinos-robustez}

Para completar a análise empírica, a Tabela~\ref{tab:h1-destinos-robustez} avalia o destino comparativo da pastagem pecuária e submete as estimativas a uma ampla bateria de testes de sensibilidade inferencial, placebo temporal e cortes amostrais.

\begin{table}[htbp]
\centering
\scriptsize
\setlength{\tabcolsep}{2.2pt}
\captionstyle{\centering}
\normalcaption
\caption{Destino comparativo pastagem, regimes temporais, placebo e sensibilidade, 2001--2024}
\label{tab:h1-destinos-robustez}
\begin{tabular}{lccccccc}
\hline
 & (1) & (2) & (3) & (4) & (5) & (6) & (7) \\
Especificação & Pastagem & 2006--2024 & Pós-2005 & Placebo $t+1$ & Clust. Duplo & Winsor 99\% & Base 3.384 \\
\hline
\multicolumn{8}{l}{\textit{Painel A: Soja ($c = s$)}} \\
$\ln P_{s,t-1} \times Z_{i,s}$ & 0,3206*** & -0,0067 & 0,0009 & 0,0559** & 0,0848* & 0,0555** & 0,0848*** \\
 & (0,1016) & (0,0083) & (0,0097) & (0,0261) & (0,0481) & (0,0230) & (0,0314) \\
$Z_{i,s} \times \mathbb{I}[t \ge 2005]$ & -- & -- & 0,0954* & -- & -- & -- & -- \\
 &  &  & (0,0514) &  &  &  &  \\
$\ln P_{s,t+1} \times Z_{i,s}$ & -- & -- & -- & 0,0844** & -- & -- & -- \\
 &  &  &  & (0,0338) &  &  &  \\
$Z_{i,s} \times \tau_t$ & -0,0485 & 0,0099* & -- & 0,0428*** & 0,0405** & 0,0359*** & 0,0456*** \\
 & (0,0369) & (0,0054) &  & (0,0153) & (0,0165) & (0,0122) & (0,0134) \\
Observações (Painel A) & 2.767 & 2.189 & 2.769 & 2.633 & 2.767 & 2.767 & 3.384 \\
$R^2$ total (Painel A) & 0,4915 & 0,5143 & 0,3770 & 0,3883 & 0,3828 & 0,4950 & 0,3711 \\
$R^2$ \textit{within} (Painel A) & 0,0195 & 0,0048 & 0,0359 & 0,0257 & 0,0198 & 0,0254 & 0,0212 \\
\hline
\multicolumn{8}{l}{\textit{Painel B: Milho ($c = m$)}} \\
$\ln P_{m,t-1} \times Z_{i,m}$ & 0,4600*** & -0,0040 & 0,0032 & 0,0198** & 0,0519* & 0,0432*** & 0,0527*** \\
 & (0,1397) & (0,0079) & (0,0076) & (0,0100) & (0,0277) & (0,0142) & (0,0146) \\
$Z_{i,m} \times \mathbb{I}[t \ge 2005]$ & -- & -- & 0,1036*** & -- & -- & -- & -- \\
 &  &  & (0,0329) &  &  &  &  \\
$\ln P_{m,t+1} \times Z_{i,m}$ & -- & -- & -- & 0,1026*** & -- & -- & -- \\
 &  &  &  & (0,0304) &  &  &  \\
$Z_{i,m} \times \tau_t$ & 0,0033 & 0,0157*** & -- & 0,0531*** & 0,0538*** & 0,0491*** & 0,0543*** \\
 & (0,0413) & (0,0060) &  & (0,0150) & (0,0192) & (0,0127) & (0,0123) \\
Observações (Painel B) & 2.767 & 2.189 & 2.769 & 2.633 & 2.767 & 2.767 & 3.384 \\
$R^2$ total (Painel B) & 0,4937 & 0,5149 & 0,3742 & 0,3867 & 0,3807 & 0,4944 & 0,3685 \\
$R^2$ \textit{within} (Painel B) & 0,0259 & 0,0083 & 0,0392 & 0,0271 & 0,0165 & 0,0246 & 0,0178 \\
\hline
\end{tabular}
\par\smallskip\begin{minipage}{\textwidth}\footnotesize
\centering
Fonte: elaboração própria com MapBiomas 10.1, Banco Mundial, BLS/FRED, FAO/IIASA (GAEZ v4) e IBGE.
\end{minipage}
\end{table}

\FloatBarrier

A coluna (1) da Tabela~\ref{tab:h1-destinos-robustez} contrasta a dinâmica das lavouras agrícolas com a conversão direta para pastagem bovina. Em consonância com o diagnóstico empírico de que mais de 90\% do desmatamento bruto de Mato Grosso converteu-se originariamente em pastagens, o gradiente de mercado na pecuária é fortemente positivo e de magnitude substancialmente mais elevada: $0{,}3206$ para soja ($p = 0{,}002$) e $0{,}4600$ para milho ($p = 0{,}001$), ambos significativos a 1\%. Como o potencial GAEZ de grãos capta atributos geofísicos de relevo plano e mecanizável, esse resultado confirma que a pecuária atuou historicamente na fronteira mato-grossense não apenas como atividade produtiva autônoma, mas como ponta-de-lança da apropriação e desmatamento de terras férteis que viriam posteriormente a ser transferidas e convertidas para o cultivo mecanizado de grãos \cite{MacedoEtAl2012,SimoesEtAl2020,JungEtAl2022}.

As colunas seguintes avaliam a robustez inferencial e os limites de identificação:
\begin{itemize}[leftmargin=1.5cm,itemsep=3pt]
    \item \textbf{Subjanela 2006--2024 (coluna 2):} No período pós-superciclo inicial e pós-PPCDAm ($N = 2.189$), os coeficientes tornam-se indistinguíveis de zero ($-0{,}0067$ para soja e $-0{,}0040$ para milho), corroborando a quebra estrutural identificada na subseção anterior;
    \item \textbf{Degrau institucional pós-2005 (coluna 3):} A substituição da tendência linear por um salto de nível na aptidão ($Z_{jm} \times \mathbb{I}[t \ge 2005]$) absorve a variação macroeconômica associada à transição de regimes, gerando coeficientes positivos no degrau ($0{,}0954$, $p = 0{,}065$ para soja; $0{,}1036$, $p = 0{,}002$ para milho);
    \item \textbf{Teste de placebo temporal (coluna 4):} A interação com preços futuros em $t+1$ gera coeficientes positivos ($0{,}0844$, $p = 0{,}014$ na soja; $0{,}1026$, $p = 0{,}001$ no milho). Esse comportamento decorre da forte autocorrelação serial dos preços de commodities em superciclos plurianuais e do caráter antecipatório das decisões de investimento em abertura de terras (cujo ciclo de destoca e gradagem demanda múltiplos anos até o primeiro plantio). Esse achado impõe prudência metodológica, indicando que as estimativas capturam tendências médias associadas aos ciclos e não choques de identificação puramente pontuais de ano para ano \cite{GoldsmithPinkhamEtAl2020};
    \item \textbf{Cluster duplo (coluna 5):} Sob agrupamento conjunto por município ($G = 137$) e por ano ($T = 24$), os coeficientes de mercado preservam significância estatística a 10\% ($p = 0{,}091$ para soja e $p = 0{,}074$ para milho), garantindo que as conclusões não decorrem de dependência espacial contemporânea residual;
    \item \textbf{Winsorização a 99\% (coluna 6):} O truncamento dos valores extremos da taxa de conversão preserva a significância a 5\% na soja ($0{,}0555$, $p = 0{,}017$) e a 1\% no milho ($0{,}0432$, $p = 0{,}003$), comprovando que os resultados não são impulsionados por \textit{outliers} isolados;
    \item \textbf{Painel balanceado expandido com zeros (coluna 7):} Na grade com imputação em $3.384$ observações, os coeficientes mantêm magnitudes comparáveis ($0{,}0848$ e $0{,}0527$) e confirmam o padrão identificado na amostra de referência.
\end{itemize}

\subsection{Diálogo com a literatura econômica e implicações para a tese}
\label{discussao-literatura-tese}

A totalidade dos achados econométricos dialoga diretamente com as três grandes vertentes da literatura de uso da terra e fronteira agrícola:
\begin{enumerate}[label=(\roman*),leftmargin=1.5cm,itemsep=3pt]
    \item \textbf{Confronto com \citeonline{BermanEtAl2023}:} Enquanto Berman et al. encontram um efeito médio positivo de preços globais sobre perda florestal agregada em nível global, os resultados deste artigo revelam que, em uma das fronteiras agrícolas mais dinâmicas do planeta, a resposta aos preços varia radicalmente conforme o uso de destino (concentrada na pecuária) e é fortemente condicionada por pactos institucionais setoriais (Moratória da Soja);
    \item \textbf{Confronto com \citeonline{HardingEtAl2021} e \citeonline{AssuncaoEtAl2015}:} Nossas estimativas estendem a compreensão da modulação institucional ao comprovar econometricamente que embargos ambientais territorializados do IBAMA amortecem em mais de 30\% a elasticidade de abertura de terras, oferecendo um freio concreto aos impulsos de mercado;
    \item \textbf{Confronto com \citeonline{MacedoEtAl2012} e \citeonline{Braganca2018}:} Comprovamos que a estabilização da conversão direta após 2005 conviveu com um expressivo ciclo de intensificação produtiva em áreas já consolidadas, viabilizado pela segunda safra de milho e pela conversão de pastagens degradadas em lavouras.
\end{enumerate}

No plano da arquitetura conceitual desta tese de doutorado, este capítulo estabelece e valida formalmente o primeiro vetor do sistema socioecológico: os incentivos econômicos de mercado operam como aceleradores da transformação da paisagem, enquanto a aptidão biofísica determina a geografia local dessa pressão e a governança institucional atua como um freio amortecedor. Ao modelar como esse arranjo interage ao longo do tempo para determinar o ritmo de conversão e a conservação da vegetação remanescente, este estudo fecha o primeiro elo analítico da tese, deixando preparado o terreno para o Capítulo 2, que investigará a via reversa: em que medida a perda acumulada desse capital natural e os impactos sobre o microclima regional afetam, retroativamente, a própria produtividade total dos fatores da agricultura mato-grossense.